\documentclass[11pt]{article}
% Shared print layout. The web's provisional design and MathJax are independent.
\usepackage[utf8]{inputenc}
\usepackage[T1]{fontenc}
\usepackage[spanish,es-nodecimaldot]{babel}
\usepackage{lmodern}
\usepackage{amsmath,amssymb}
\usepackage{graphicx}
\usepackage{geometry}
\geometry{a4paper,left=20mm,right=20mm,top=25mm,bottom=22mm,headheight=16pt,headsep=8mm}
\usepackage{fancyhdr,lastpage}
\usepackage{enumitem}
\usepackage{xcolor}
\usepackage{pgfplots}
\pgfplotsset{compat=1.18}
\usepgfplotslibrary{fillbetween}
\usepackage[hidelinks,pdfencoding=auto]{hyperref}
\hypersetup{pdflang=es-ES,pdfauthor={},pdfsubject={Matemáticas Aplicadas a las Ciencias Sociales II}}
\pagestyle{fancy}
\fancyhf{}
\fancyhead[L]{\small Matemáticas Aplicadas a las Ciencias Sociales II}
\fancyhead[R]{\small 2.\textsuperscript{o} Bachillerato}
\fancyfoot[L]{\footnotesize Recuperación y repaso | Edición revisada 1.0}
\fancyfoot[R]{\footnotesize Página \thepage\ de \pageref{LastPage}}
\renewcommand{\headrulewidth}{0.4pt}
\setlength{\parindent}{0pt}
\setlength{\parskip}{0.5em}
\setlist{itemsep=0.3em,topsep=0.3em}
\newcommand{\apartado}[1]{\subsection*{#1}}
\pgfplotsset{
  revision/.style={
    width=0.94\linewidth,height=7.5cm,
    axis lines=middle,grid=major,grid style={gray!18},
    tick label style={font=\small},label style={font=\small},
    xlabel={$x$},ylabel={$y$},samples=160,
    legend style={at={(0.5,-0.16)},anchor=north,draw=none,font=\small},
    clip=true
  }
}

\hypersetup{pdftitle={Recuperación de Inferencia y repaso - Examen revisado}}
\begin{document}
\begin{minipage}{0.17\linewidth}\includegraphics[width=0.8\linewidth]{logo-al-andalus.png}\end{minipage}\hfill
\begin{minipage}{0.79\linewidth}{\Large\bfseries Recuperación y repaso}\\[0.4em]
Duración original: \textbf{55 minutos}.\end{minipage}

Apellidos y nombre: \hrulefill

Fecha: \hrulefill

Se permite calculadora. Responde razonadamente y numera tus páginas. Cada ejercicio incluye 0,5 puntos de expresión, sin añadirlos a su máximo.

\textbf{Nota de edición:} se conservan los nueve ejercicios de 2,5 puntos (22,5 puntos nominales). El original no especifica selección ni conversión de nota; el docente debe aclarar estas condiciones antes de utilizarlo como examen. No se impone aquí una regla nueva.



\section*{Ejercicio 1 \hfill 2,5 puntos}  (Distribuciones) Se ha tomado una muestra de 16 pacientes tratados por un especialista y se ha observado que el tiempo de espera en su consulta, en minutos, ha sido de:

\begin{center}
8,\ 9.2,\ 10,\ 8.5,\ 12,\ 9,\ 11.3,\ 7,\ 8.5,\ 8.3,\ 7.6,\ 9,\ 9.4,\ 10.5,\ 8.9,\ 6.8
\end{center}

Supóngase que el tiempo de espera se distribuye según una ley Normal de varianza 4 y media desconocida.

\begin{enumerate}[label=(\alph*)]

  \item (1,5 puntos)  Halle un intervalo de confianza al 97.5\% para estimar el tiempo medio de espera de los pacientes tratados por este especialista.

  

  \item (1 punto)  \textquestiondown{}Cuál debería ser el tamaño mínimo de la muestra para asegurar, con un nivel de confianza del 90\%, que el error cometido sea, a lo sumo, de 0.3 minutos?

  

\end{enumerate}



\section*{Ejercicio 2 \hfill 2,5 puntos} (Distribuciones) El precio de venta al público del kilogramo de frambuesas sigue una ley Normal de media desconocida y varianza 9. En una localidad se eligen 10 comercios de manera aleatoria, obteniéndose los siguientes precios en euros:

\begin{center}
12.3,\ 10,\ 9.1,\ 11,\ 10.5,\ 11.8,\ 9.9,\ 11.5,\ 10.9,\ 13
\end{center}

\begin{enumerate}[label=(\alph*)]

  \item (0,5 puntos)  \textquestiondown{}Qué distribución siguen las medias de las muestras de tamaño 10?

  

  \item (1 punto)  Con los datos obtenidos de la muestra, determine un intervalo de confianza al 97\% para el precio medio del kilogramo de frambuesas.

  

  \item (1 punto)  Con el mismo nivel de confianza, calcule el tamaño mínimo que debe tener una muestra para que el error cometido al estimar el precio medio del kilogramo de frambuesas sea menor a 1.5 euros.

  

\end{enumerate}


\clearpage
\section*{Ejercicio 3 \hfill 2,5 puntos}  (Distribuciones) Se sabe que la longitud, en centímetros, de una especie de estrella de mar en una determinada zona sigue una ley Normal con desviación típica 3. Para estimar la longitud media, se extrae una muestra de tamaño 36 y se obtiene el intervalo de confianza \( (6.04,\ 8) \) al 95\%. Se pide:

\begin{enumerate}[label=(\alph*)]

  \item (0,5 puntos)  Calcule la media muestral.

  

  \item (0,5 puntos)  Calcule el error de estimación máximo cometido.

  

  \item (1 punto)  Si aumentamos el tamaño muestral a 49, \textquestiondown{}qué efecto produce sobre el error máximo cometido? Calcule este error.

  

  \item (0,5 puntos)  Si aumentamos el nivel de confianza, \textquestiondown{}qué efecto produce sobre el error de estimación máximo? Justifique la respuesta.

  

\end{enumerate}


\section*{Ejercicio 4 \hfill 2,5 puntos} (Probabilidad) Se han mezclado 90 llaves electrónicas de apertura de un determinado garaje, con apariencia idéntica, de las cuales 60 funcionan correctamente y 30 no funcionan. Se eligen al azar 2 de las 90 llaves.

\begin{enumerate}[label=(\alph*)]

  \item (0,7 puntos)  \textquestiondown{}Cuál es la probabilidad de que las dos llaves elegidas abran la puerta del garaje?

  

  \item (0,8 puntos)  \textquestiondown{}Cuál es la probabilidad de poder abrir el garaje con alguna de ellas?

  

  \item (1 punto)  \textquestiondown{}Cuál es la probabilidad de que una de las llaves elegidas funcione correctamente y la otra no?

  

\end{enumerate}



\clearpage
\section*{Ejercicio 5 \hfill 2,5 puntos} (Probabilidad) Una empresa almacena el mismo número de latas de refresco de cola, naranja y limón. De las 30.000 latas de refresco almacenadas, se sabe que 1800 latas de cola, 2400 de naranja y 3000 de limón caducan en 2021.

\begin{enumerate}[label=(\alph*)]

  \item (1,5 puntos)  \textquestiondown{}Cuál es la probabilidad de que una lata elegida al azar caduque en 2021?

  

  \item (1 punto)  Si se ha elegido al azar una lata que no caduca en 2021, \textquestiondown{}cuál es la probabilidad de que sea de cola?

  

\end{enumerate}






\section*{Ejercicio 6 \hfill 2,5 puntos} (SELs y ProgLin)

\begin{enumerate}[label=(\alph*)]

  \item (1 punto)  Una fábrica de electrodomésticos dispone de dos cadenas de montaje. En una hora de trabajo, la cadena A produce 10 lavadoras y 5 frigoríficos, mientras que la cadena B produce 7 lavadoras y 6 frigoríficos. El coste por hora de trabajo es de 1200 euros para A y 1500 euros para B. La cadena A puede trabajar, como máximo, el doble de horas que la cadena B. Si se deben producir al menos 400 lavadoras y 280 frigoríficos, formule, \textbf{sin resolver}, el problema que permite obtener las horas de funcionamiento de A y B para minimizar el coste.

  

  \item (1,5 puntos)  Represente el recinto definido por las siguientes inecuaciones y calcule sus vértices:
  \[
  x + 2y \geq 7,\quad 4x - y \geq 1,\quad 2x - y \leq 4,\quad 3x + 2y \leq 20,\quad x \geq 0,\quad y \geq 0
  \]
  Obtenga el valor mínimo de la función \( F(x, y) = 2x + y \) en el recinto anterior, así como el punto en el que se alcanza.

  

\end{enumerate}




\clearpage
\section*{Ejercicio 7 \hfill 2,5 puntos} (Análisis) Se considera la función \( f(x) = ax^3 + bx + 4 \), con \( a \) y \( b \) números reales.

\begin{enumerate}[label=(\alph*)]

  \item (1 punto)  Determine los valores \( a \) y \( b \) para que \( f \) tenga un extremo relativo en el punto \( (2,\ 36) \).

  

  \item (0,75 puntos)  Para \( a = 4 \) y \( b = -3 \), estudie la monotonía de \( f \) y determine sus extremos relativos.

  

  \item (0,75 puntos)  Para \( a = 4 \) y \( b = -3 \), calcule la función \( F(x) \) que verifica \( F'(x) = f(x) \) y \( F(2) = 10 \).

  

\end{enumerate}



\section*{Ejercicio 8 \hfill 2,5 puntos} (Derivadas e integrales)

\begin{enumerate}[label=(\alph*)]

  \item (1,2 puntos)  Calcule la derivada de las siguientes funciones:
  \[
  f(x) = (-5 + x^2)^2 \cdot e^{3x} \quad\quad 
  g(x) = \frac{\ln(x^3 - 5x)}{1 - x^2}
  \]

  

  \item (1,3 puntos)  Calcule el área del recinto acotado por la gráfica de \( h(x) = -x^2 + 2x + 3 \) y el eje de abscisas.

  

\end{enumerate}



\clearpage
\section*{Ejercicio 9 \hfill 2,5 puntos} (Matrices) Sean \( A,\ B,\ X,\ Y \) matrices invertibles que verifican:
\[
A \cdot X = B \quad \text{y} \quad B \cdot Y = A
\]

\begin{enumerate}[label=(\alph*)]

  \item (1 punto)  Compruebe que \( Y^{-1} = X \).

  

  \item (1,5 puntos)  Para \( A = \begin{pmatrix} 1 & 2 \\ 1 & 3 \end{pmatrix} \) y 
  \( B = \begin{pmatrix} 2 & 1 \\ 0 & -1 \end{pmatrix} \), halle \( X \) e \( Y \).

  

\end{enumerate}






\end{document}
